\section{Results}

\subsection{RQ1 and RQ2: knowledge sources on Dart}

Table~\ref{tab:dart} gives the divergence rate of every run.
Human knowledge reproduces our previous result on fresh seeds (30.9\%,
against 28.8\% before). Code-read knowledge sits between human and none, and
every run in each of these three arms is ordered: all five human runs exceed
all five code runs, which exceed all five \emph{none} runs.

\begin{table}[t]
\centering
\caption{Dart divergence rate per run (\% of all generated documents),
seeds 700--704 in order.}
\label{tab:dart}
\begin{tabular}{lrl}
\toprule
Arm & Mean & Per seed \\
\midrule
none  & 1.45  & 3.40, 0.00, 0.67, 3.20, 0.00 \\
human & 30.89 & 21.60, 29.65, 48.30, 29.13, 25.75 \\
code  & 11.55 & 10.10, 16.04, 8.40, 13.40, 9.80 \\
probe & 10.51 & 1.07, \textbf{49.27}, 1.27, 0.93, 0.00 \\
\bottomrule
\end{tabular}
\end{table}

\begin{table}[t]
\centering
\caption{Pairwise Dart comparisons: difference in mean rate (percentage
points), exact two-sided Mann--Whitney U, Holm-adjusted $p$, Cliff's
$\delta$. $n{=}5$ per arm.}
\label{tab:dart-tests}
\begin{tabular}{lrrrr}
\toprule
Comparison & $\Delta$ & $p$ & $p_{\mathrm{Holm}}$ & $\delta$ \\
\midrule
probe$-$none (primary) & $+9.05$  & 0.690 & 0.690 & $+0.20$ \\
probe$-$human (co-prim.) & $-20.38$ & 0.151 & 0.452 & $-0.60$ \\
human$-$none & $+29.43$ & 0.008 & 0.048 & $+1.00$ \\
code$-$none & $+10.09$ & 0.008 & 0.048 & $+1.00$ \\
code$-$human & $-19.34$ & 0.008 & 0.048 & $-1.00$ \\
probe$-$code & $-1.04$  & 0.151 & 0.452 & $-0.60$ \\
\bottomrule
\end{tabular}
\end{table}

\textbf{RQ1.} Probe knowledge is not detectably better than none
(Table~\ref{tab:dart-tests}, $p \approx 0.69$, $\delta = +0.20$). By the
pre-registered rule this is ``not detected at $n{=}5$'', and the claim of
comparability to human knowledge is not supported ($\delta = -0.60$).
The mean hides the shape of the distribution. Four probe runs stay at or
below 1.3\%, in the range of \emph{none}; one, seed 701, reaches 49.3\%,
above every human run. Seed 701 is the only seed whose probing phase
observed a divergence at all: one of its 100 probes across five seeds
diverged (\texttt{built\_value} accepting a nested \texttt{tags: null}), and
its summary reported it. The probes of the other four seeds diverged on
none of their 20 documents. When probing sees nothing, its summary is a
confident description of agreement, and the loop behaves as if it had no
knowledge. Summaries can also overgeneralize: one states that no
implementation throws an undocumented exception, although three of the four
paths throw Dart's private \texttt{\_TypeError} on the probes it ran.

\textbf{RQ2.} Human $>$ code $>$ none holds with the largest possible
effect in every pair and survives the Holm correction. Code-read knowledge is
the stable middle: it varies between 8.4\% and 16.0\% and never fails.
Probe knowledge has a similar mean but the opposite profile. The two are not
distinguishable at $n{=}5$.

\begin{table}[t]
\centering
\caption{Dart pattern recall. \emph{Stated}: summaries (of five) that state
the fact behind the pattern. \emph{Found}: runs (of five) that found the
pattern at least once.}
\label{tab:recall}
\begin{tabular}{lcccc}
\toprule
& \multicolumn{2}{c}{\texttt{RAAR} (\texttt{id})} & \multicolumn{2}{c}{\texttt{RRRA} (\texttt{tags})} \\
Arm & stated & found & stated & found \\
\midrule
none  & --   & 3/5 & --   & 3/5 \\
human & 5/5  & 5/5 & 5/5  & 5/5 \\
code  & 5/5  & 5/5 & 0/5  & 1/5 \\
probe & 0/5  & 0/5 & 1/5$^{*}$ & 4/5 \\
\bottomrule
\end{tabular}\\[2pt]
{\footnotesize $^{*}$Seed 701 states that \texttt{built\_value} accepts
\texttt{tags: null}: related, but not the missing-field fact.}
\end{table}

\textbf{What the knowledge contains decides what is found.}
Table~\ref{tab:recall} compares, for each Dart pattern, whether an arm's
knowledge states the fact behind it with how many runs found it. Every
code summary states that two paths accept a non-integer \texttt{id} through
\texttt{num.toInt()}; every code run finds \texttt{RAAR}. None of them
mentions large integers or overflow, which is the other half of the
mechanism, but a generator that varies \texttt{id} types reaches it anyway.
No code summary states that \texttt{built\_value} substitutes an empty
list for a missing \texttt{tags}, and one code run finds \texttt{RRRA}.
The probe summaries are the mirror image: none mentions a non-integer
\texttt{id}, and no probe run finds \texttt{RAAR}, not even the one at
49\%, whose divergences are all \texttt{RRRA}. No arm found a Dart pattern
outside the answer key. The loop exploits what it is told; it does not
extend it, which is what our previous study observed for human knowledge.

\textbf{Cost.} A \emph{none} or \emph{human} run takes five calls, about
\$0.014 at list prices. Code knowledge adds two calls per seed (\$0.004) and
probing four (\$0.006), plus 20 harness executions. Knowledge acquisition is
cheap compared with the loop; the whole study, including pilots, cost less
than one US dollar.

\subsection{RQ3: Kotlin}

On Kotlin the picture is reversed. The probing agents observed divergence
on 11 to 14 of their 20 probes on every seed: the four libraries disagree
about wrong types, nulls, unknown keys, unknown enum constants and lenient
syntax. The summaries record these disagreements, and the loop uses them
(Table~\ref{tab:kotlin}).

\begin{table}[t]
\centering
\caption{Kotlin divergence rate (\% of all documents), probe vs.\ none.
The replication on seeds 805--809 was declared before it ran.}
\label{tab:kotlin}
\begin{tabular}{lrrrr}
\toprule
Seeds & none & probe & $p$ & $\delta$ \\
\midrule
800--804 (pre-registered) & 45.77 & 64.34 & 0.095 & $+0.68$ \\
805--809 (replication)    & 54.14 & 69.92 & 0.056 & $+0.76$ \\
pooled, $n{=}10$ per arm  & 49.95 & 67.13 & \textbf{0.007} & $+0.70$ \\
\midrule
pooled, RFC 8259-valid only & 42.55 & 56.49 & 0.075 & $+0.48$ \\
\bottomrule
\end{tabular}
\end{table}

The pre-registered comparison on seeds 800--804 had a large effect but did
not reach significance ($p \approx 0.095$). As the design allowed, we
declared a replication on fresh seeds before running it; it showed the same
direction and size ($+15.8$ points, $\delta = +0.76$), and the pooled
comparison over ten seeds per arm is significant ($p \approx 0.007$). We
report all three rows: neither half is significant alone.

The last row restricts both arms to documents that are valid JSON under
RFC~8259~\cite{rfc8259}. The harness's own validity label turned out to be
too lenient during triage (Section~\ref{sec:triage}), so this row uses a
strict parser applied after the fact, as recorded in the deviations log.
The effect shrinks ($\delta = +0.48$, $p \approx 0.075$): part of the
probe arm's advantage comes from malformed input that some libraries accept
and others do not.

Every Kotlin run in both arms, including all ten without knowledge,
rediscovered the Gson null-safety bypass, in 114 to 970 documents per run.
On Kotlin this answer key is too easy to separate the arms.

\subsection{Triage of Kotlin divergences}
\label{sec:triage}

We grouped every divergent Kotlin document by pattern and exception types,
reproduced each group by hand with a minimal document that changes one
field of a valid record, and checked the result against each library's
documentation and issue tracker. No run produced a value divergence:
whenever all four libraries accepted a document, they decoded the same
value.

Eleven behaviours are documented defaults or known limitations. They are
Gson's null-safety bypass and its mapping of unknown enum constants to
\texttt{null}; Gson and Jackson decoding a missing or \texttt{null}
\texttt{id} as 0, which Jackson's Kotlin module documents as a caveat;
Gson, Moshi and Jackson accepting \texttt{null} elements in a
\texttt{List<String>} (strict null checks are opt-in for Jackson); the coercion
of numbers and booleans to strings; kotlinx.serialization and Jackson rejecting
unknown keys; kotlinx.serialization requiring a nullable property with no
default; Jackson truncating \texttt{1.9} to \texttt{1}; and Gson's lenient
syntax; and two coercions of string values to \texttt{Long}: Jackson reading
an empty string as 0, and Gson and Moshi accepting a numeric string with a
fraction, such as the string \texttt{1.0}.

One behaviour is a standards violation. kotlinx.serialization's
\emph{default}, non-lenient \texttt{Json} accepts raw control characters
(U+0000--U+001F) inside strings, which RFC~8259 forbids and which Jackson
and Dart's \texttt{jsonDecode} reject. We reproduced it outside the harness
on version 1.11.0 and on the latest release candidate, 1.12.0-RC. It had
been reported only inside a generic benchmark issue that was closed as not
planned without comment~\cite{kotlinxissue3273}, and we have prepared a
focused report. The same leniency is why the harness's validity label,
which used kotlinx.serialization's parser, was not RFC-strict.

Two more defects surfaced only in our manual reproductions, not in any run
or probe: Gson saturating an integer literal just above $2^{63}$ to
\texttt{Long.MAX\_VALUE}~\cite{gsonissue1727}, and kotlinx.serialization
wrapping integers modulo $2^{64}$, so that \texttt{18446744073709551617}
decodes as \texttt{1}~\cite{kotlinxissue3267}. Both are known and open. We
do not credit them to the method: like \texttt{RAAR} for the probe arm,
they sit in a narrow numeric region that neither the probes nor the
generators visited.
